\documentclass{dsekprotokoll}

\usepackage[T1]{fontenc}
\usepackage[utf8]{inputenc}
\usepackage[swedish]{babel}
\usepackage{tikz}
\usetikzlibrary{arrows.meta,positioning}

\setheader{Policy för Hantering av nya poster}{Policydokument}{}

\title{Policy för Hantering av nya poster}
\author{D-sektionen}

\begin{document}
\maketitle

\section{Formalia}

\subsection{Sammanfattning}
Denna policy beskriver hur D-sektionen inom TLTH arbetar med frågor som rör
hanteringen och införandet av nya funktionärsposter på D-sektionen.

\subsection{Syfte}
Syftet med denna policy är att skapa ett tydligt och enhetligt tillvägagångssätt
för hur nya poster inom D-sektionen kan införas, namnges och beslutas om.
Policyn säkerställer att posterna blir en naturlig och transparent del av
sektionens organisation, att styrdokument och verklighet överensstämmer samt att
sektionens verksamhet framstår som tydlig och tillgänglig för både medlemmar och
utomstående.

\subsection{Omfattning}
Policyn gäller för alla ännu ej befintliga funktionärsposter som kan komma att
uppstå inom sektionen, oavsett om de är kopplade till ett utskott, till en
särskild verksamhet eller avser en mer fristående funktion. Policyn omfattar
inte redan befintliga poster som regleras i sektionens reglemente.

\subsection{Historik}
Utkast färdigställt av: Måns Bard Nilsson.\\
Ursprungligen antagen enligt beslut: HTM1 2025.

\section{Tillvägagångssätt}

\subsection{Förslag om ny post}
När en ny post föreslås ska detta dokumenteras skriftligt och innehålla en
motivering till varför posten behövs, vilka uppgifter som förväntas ingå samt
hur posten förhåller sig till sektionens övriga organisation. Denna
dokumentation ska tillhandahållas som underlag inför ett beslutande organ.

Detta ska göras innan det att en person kan bli vald till en post då posten inte
existerar före dess genomröstande.

\subsection{Beslutande organ}
Beslut om att inrätta en ny post fattas av sektionsmötet. Förslaget kan väckas
av styrelsen, ett utskott eller enskilda medlemmar genom motion eller
proposition.

\subsection{Benämning och titel}
Postens namn ska vara tydligt, lätt att förstå och konsekvent användas i alla
sammanhang (hemsida, protokoll, utskottsstruktur).

Namnet ska inte riskera att förväxlas med redan befintliga poster eller externa
roller.

\subsection{Reglemente och styrdokument}
I samband med tillhandahållandet av förslag för en ny post ska även resterande
styrdokument med korrekta ändringar tillhandahållas som bilagor för beslutet.

Posten blir del av sektionens förteckning av funktionärer då ett beslut att
införa en ny funktionärspost är justerat.

\subsection{Översyn och avveckling}
Nya poster ska följas upp av styrelsen i samråd med utskottsordförande för
utskottet som posten befinner sig inom under det kommande verksamhetsåret. Om
posten inte fyller en funktion eller fungerar väl i praktiken kan detta lyftas
på ett kommande sektionsmöte som kan komma att avveckla den berörda posten.

\section{Tillämpning}
Det är sektionsstyrelsens ansvar att se till att denna policy följs. Policyn ska
verka som ett komplement till våra övriga styrdokument.

\end{document}
